\thispagestyle{plain}
\begin{center}
  \section*{Abstract}
\end{center}
\begingroup
This study investigates the role of an aircraft's wing sweep angle as a design variable in aerodynamic shape optimization (ASO) for conventional aircraft configurations, focusing on the NASA Common Research Model (CRM) and A320neo wings. While ASO has advanced significantly in recent decades, the wing sweep angle remains under-explored despite its significance in influencing wave drag at transonic speeds. By conducting ASO at discrete and variable sweep angles, this study reveals that both wings exhibit optimal aerodynamic performance between 45-50 degrees of wing sweep, with notable drag reduction compared to their original geometries. However, the benefits diminish beyond the nominal sweep angle as a result of weight and induced drag penalties, highlighting the complex interplay between the wing sweep angle and overall aircraft drag. These findings underscore the need for comprehensive consideration of design variable interactions in aircraft design optimization.
\par
\endgroup
\cleardoublepage